% =====================================================================
%  PROPOSAL P3 -- The certificate as a compliance artifact
%
%  Standalone. Set this as Overleaf's main document and compile with
%  pdfLaTeX; Overleaf runs BibTeX automatically. Locally:
%    pdflatex P3 && bibtex P3 && pdflatex P3 && pdflatex P3
% =====================================================================
\documentclass[11pt,a4paper]{article}
% =====================================================================
%  Shared preamble for the three research proposals.
%
%  Each proposal \input{}s this file and is otherwise standalone, so any
%  one of them can be set as Overleaf's main document and compiled on its
%  own. Nothing here is specific to a single proposal.
% =====================================================================
\usepackage[T1]{fontenc}
\usepackage[utf8]{inputenc}
% Times, matching the parent manuscript. mathptmx is in every TeX Live and on
% Overleaf; lmodern is deliberately not used, since it is absent from minimal
% distributions and this preamble should compile anywhere.
\usepackage{mathptmx}
\usepackage[a4paper,margin=25mm]{geometry}
\usepackage{amsmath,amssymb,amsthm}
\usepackage{booktabs}
\usepackage{tabularx}
\usepackage{xcolor}
\usepackage{titlesec}
\usepackage{enumitem}
\usepackage{microtype}
\usepackage[hidelinks,breaklinks]{hyperref}

% ---- shared style tokens, matching the parent manuscript ------------
\definecolor{netcol}{RGB}{31,78,121}
\definecolor{autocol}{RGB}{18,120,90}
\definecolor{bridgecol}{RGB}{176,96,0}
\definecolor{accent}{RGB}{140,20,20}

\newcommand{\layernet}{\textcolor{netcol}{\textbf{network}}}
\newcommand{\layerauto}{\textcolor{autocol}{\textbf{autonomy}}}
\newcommand{\dataset}[1]{\textsc{#1}}
\newcommand{\code}[1]{\texttt{#1}}
\newcommand{\Real}{\mathbb{R}}
\newcommand{\Norm}[1]{\left\lVert #1\right\rVert}
\newcommand{\MCR}{\mathrm{MCR}}

\newtheorem{claim}{Claim}
\newtheorem{proposition}{Proposition}
\newtheorem{definition}{Definition}
\theoremstyle{remark}
\newtheorem{remark}{Remark}

% ---- section styling -------------------------------------------------
\titleformat{\section}{\large\bfseries\color{netcol}}{\thesection}{0.7em}{}
\titleformat{\subsection}{\normalsize\bfseries}{\thesubsection}{0.6em}{}
\titlespacing*{\section}{0pt}{1.4em}{0.6em}

\setlist[itemize]{leftmargin=1.35em,topsep=0.3em,itemsep=0.25em}
\setlist[enumerate]{leftmargin=1.6em,topsep=0.3em,itemsep=0.25em}

% ---- a boxed "what this rests on" callout ---------------------------
\newenvironment{honestly}
  {\par\medskip\noindent\begin{tabular}{@{}p{2.5mm}@{\hspace{3mm}}p{0.9\linewidth}@{}}
   \textcolor{accent}{\rule{2.5mm}{\baselineskip}} &
   \small\itshape}
  {\end{tabular}\par\medskip}

% ---- the proposal header --------------------------------------------
\newcommand{\proposalheader}[3]{%
  \begin{center}
  {\Large\bfseries #1\par}\medskip
  {\large\itshape #2\par}\medskip
  {\small #3\par}
  \end{center}
  \medskip\hrule\medskip}


\begin{document}

\proposalheader
  {P3 --- The Certificate as a Compliance Artifact}
  {Certified Robustness Meets Airworthiness: Mapping ML Robustness Guarantees
   onto SORA Operational Volumes and ED-324 Learning Assurance}
  {Research proposal, RobustUAVs.ai project $\cdot$ Roger Nick Anaedevha $\cdot$
   drafted 2026-08-13\\
   Follow-on to \emph{From Wire to Flight} (NDSS submission)}

\begin{abstract}
\noindent
The parent manuscript already computes quantities that aviation regulators
independently require, and does not say so loudly enough. Its corridor-margin
family was chosen to bracket the JARUS SORA contingency volume; its certified
tube radius is directly comparable to that volume's lateral extent; and the
``operating region'' qualifier on its Lipschitz hypothesis is, in the vocabulary
of EASA's learning-assurance framework, an Operational Design Domain. This
proposal formalises the correspondence, performs a gap analysis against the
ED-324 learning-assurance objectives, and delivers an evidence generator that
emits a SORA-shaped compliance pack from a certified run. It is a
systematisation-and-tooling contribution, and is proposed as such.
\end{abstract}

% =====================================================================
\section{The gap}
% =====================================================================
Three developments have converged while the certified-robustness literature has
looked the other way.

\paragraph{The operational-volume methodology is settled.}
Under the JARUS Specific Operations Risk Assessment, the operational volume
comprises the flight geography plus a contingency volume whose minimum lateral
extent is the sum of GNSS position accuracy, position-holding error, map error,
a reaction-distance term, and the stopping distance of the contingency
manoeuvre~\cite{jarus_sora25}. The parent manuscript's margin family
$\{2,5,10,20\}$\,m was chosen to bracket exactly this, and its default values
for the first three terms already total about $7$\,m.

\paragraph{The AI assurance framework is arriving.}
ED-324 / ARP6983 is EUROCAE WG-114 and SAE G-34's process standard for
developing and certifying aeronautical products that implement
AI~\cite{ed324,eurocae_wg114}. EASA's learning-assurance framework turns on a
Concept of Operations, an Operational Domain, and an Operational Design Domain
for the machine-learning constituent~\cite{easa_mleap,easa_ai_concept2}, with
the broader trajectory set out in the agency's AI
roadmap~\cite{easa_ai_roadmap2}.

\paragraph{The airspace framework is becoming mandatory.}
Network Remote ID is becoming a condition of participation in
U-space~\cite{easa_uspace2021,easa_remoteid_2026}, which makes the connectivity
assumptions of a cross-layer security argument regulatory rather than
architectural.

Meanwhile the parent manuscript states its Lipschitz hypothesis as holding
``over the operating region'', and then enumerates that region: attitudes away
from gimbal-lock singularities, airspeeds below stall, actuators off their
saturation limits, no contact dynamics. \textbf{That is an ODD.} It is described
in exactly those terms and the connection is never claimed.

% =====================================================================
\section{The claim}
% =====================================================================
A three-part contribution: a mapping, a gap analysis, and a generator.

\subsection{The mapping}
A correspondence between certified-robustness vocabulary --- certified radius,
trajectory tube, operating region, certified window, mission floor --- and
SORA/ED-324 vocabulary --- contingency volume, ODD, operational domain, learning
assurance objective. Where the correspondence is exact, prove it. Where it is
not, name what is missing.

\begin{honestly}
The most likely honest finding is a negative one: \textbf{certified radii are
per-input, while assurance objectives are per-operation}. A statement that a
model is stable for every perturbation of norm at most $R$ does not by itself
discharge an objective phrased over a flight. The parent manuscript's
conjunction of a spatial and a temporal predicate is a step across that divide,
which is precisely why it is worth writing this paper --- but the gap should be
reported, not smoothed over. A systematisation whose finding is ``it all lines
up'' would be worth very little.
\end{honestly}

\subsection{The gap analysis}
For each ED-324 learning-assurance objective, classify what a
Gr\"onwall~\cite{gronwall1919}, randomized-smoothing~\cite{cohen2019smoothing},
PAC-Bayes~\cite{mcallester1999pac} or regret-based~\cite{arora2012mwu}
certificate can do:

\begin{itemize}
\item \textbf{Discharge:} the certificate is sufficient evidence for the
objective.
\item \textbf{Evidence:} the certificate contributes but does not suffice.
\item \textbf{Untouched:} the objective is about something certificates do not
address --- data governance, requirements traceability, human oversight.
\end{itemize}

The runtime-assurance reading already present in the parent manuscript ---
Simplex-style handover when the aircraft leaves the linearisable
envelope~\cite{sha2001simplex} --- is the natural bridge for objectives the
certificate cannot discharge alone, and should be developed here rather than
left as a remark.

\subsection{The evidence generator}
Emit a SORA Annex~A-shaped evidence pack from a certified run: the ODD as the
measured operating region, the contingency extent as the certified tube radius,
the mission floor as the safety objective, and every figure carrying its
provenance tag. This is what makes the paper more than a survey, and it is
directly executable against the parent project's committed results.

% =====================================================================
\section{What is reused, what is new}
% =====================================================================
\begin{center}
\small
\begin{tabularx}{\linewidth}{@{}XX@{}}
\toprule
\textbf{Reused unchanged} & \textbf{New} \\
\midrule
the entire compute side: engine, corridor family, provenance ledger &
the vocabulary mapping and its proofs \\
the four-certificate architecture &
the objective-by-objective gap analysis \\
the results already committed &
the evidence generator over \code{results/} \\
\bottomrule
\end{tabularx}
\end{center}

\noindent
This is the cheapest of the three proposals to execute and the only one with no
measurement risk.

% =====================================================================
\section{Relation to the parent manuscript's limitations}
% =====================================================================
P3 closes none of the manuscript's stated limitations, and should not be sold as
if it did. Its value is that it changes who the result is \emph{for}: from a
security audience that reads a certified floor as a bound, to a certification
audience that reads it as evidence. That is a different contribution, not a
smaller one, but the distinction should be explicit in the framing.

% =====================================================================
\section{Platform integration}
% =====================================================================
A \textbf{Compliance Pack} export: one action producing the evidence bundle for
a chosen operating point, slotting into the existing export menu alongside the
JSON, CSV, PNG and PDF exporters.

% =====================================================================
\section{Risks}
% =====================================================================
\begin{itemize}
\item \textbf{It is a systematisation paper and must be honest about that.} Its
contribution is the mapping and the generator, not new theory. Pitch it to a
venue that wants that, rather than overselling it at a top-tier security
conference and being rejected for the right reason.
\item \textbf{The standards move.} ED-324 is in preparation and the
bibliography already records it as such~\cite{ed324}; the Commission's drone
security package is in progress. Any claim keyed to a draft clause number will
rot. Key claims to concepts --- ODD, operational domain, contingency volume ---
which are stable, and cite clause numbers only in an appendix.
\item \textbf{Regulatory over-claim.} Nothing in this paper certifies an
aircraft. It maps a mathematical guarantee onto an evidence framework. The
distinction must survive the abstract.
\end{itemize}

% =====================================================================
\section{Venue and timeline}
% =====================================================================
Target an IEEE S\&P SoK track, ACM Computing Surveys, or the Digital Avionics
Systems Conference --- the last of which would put the work in front of the
people who actually write these standards. About three months, and it can run in
the gaps of P1 and P2.

\small
\bibliographystyle{plain}
\bibliography{proposals}

\end{document}
